\chapter{Implementation}
\label{ch:implementation}

This chapter describes the software implementation of the RCC agent-based model, covering the selection of the Repast4Py framework, the overall architecture, the custom spatial index design, agent serialization conventions, the porting process from Mesa, the glucose field implementation, and the simulation entry point.

% ======================================================================
\section{Framework Selection}
\label{sec:framework-selection}

The choice of simulation framework is fundamental to the scalability and maintainability of any agent-based model. After evaluating the two most prominent Python-based ABM frameworks---Mesa~\cite{mesa} and Repast4Py~\cite{repast4py}---the latter was selected for the RCC tumor microenvironment model.

\subsection{Rationale for Repast4Py}

Repast4Py offers several features that are critical for the present model:

\begin{enumerate}[label=(\roman*)]
    \item \textbf{Native 3D grids.} Repast4Py's \texttt{SharedGrid} class supports arbitrary dimensionality out of the box. In Mesa, only 2D grids are provided natively; extending them to three dimensions requires custom subclassing of \texttt{MultiGrid} with non-trivial modifications to the neighbor-query logic.

    \item \textbf{MPI parallelism via mpi4py.} Repast4Py is designed around the Message Passing Interface (MPI) through the \texttt{mpi4py} library~\cite{mpi4py}. Agents are managed within a \texttt{SharedContext} that can distribute them across MPI ranks. This architecture enables future scaling to large grids without architectural changes to the model code.

    \item \textbf{SharedContext for distributed agents.} The \texttt{SharedContext} abstraction provides a unified agent container that handles ghost-agent synchronization across rank boundaries, which is essential for multi-rank parallel execution.

    \item \textbf{Python ecosystem compatibility.} Repast4Py is a pure Python library (with \texttt{mpi4py} as its primary compiled dependency), and integrates seamlessly with NumPy, SciPy, Optuna, Pandas, and Matplotlib---all of which are used in the RCC model.
\end{enumerate}

\subsection{Comparison with Mesa}

Table~\ref{tab:mesa-vs-repast4py} summarizes the key differences between Mesa and Repast4Py that influenced the framework selection.

\begin{table}[H]
\centering
\caption{Feature comparison between Mesa and Repast4Py for the RCC model.}
\label{tab:mesa-vs-repast4py}
\begin{tabular}{@{}lll@{}}
\toprule
\textbf{Feature} & \textbf{Mesa} & \textbf{Repast4Py} \\
\midrule
Grid & \texttt{MultiGrid} (custom 3D) & \texttt{SharedGrid} (native 3D) \\
Parallelism & None (single process) & MPI via \texttt{mpi4py} \\
Scheduling & Mesa scheduler & \texttt{schedule.init\_schedule\_runner} \\
Agent base class & \texttt{mesa.Agent(model)} & \texttt{repast4py.core.Agent(id, type, rank)} \\
Border handling & Torus / Non-torus & Sticky / Periodic \\
Occupancy & Single / Multi & Single / Multiple \\
\bottomrule
\end{tabular}
\end{table}

The most consequential difference for this project is the parallelism model. Mesa is inherently single-process: all agents live in a single Python interpreter, and there is no built-in mechanism for distributing computation across cores or nodes. Repast4Py, by contrast, is architected from the ground up around MPI, making it possible to scale the simulation to larger tissue volumes in the future without redesigning the agent management infrastructure.

% ======================================================================
\section{Architecture}
\label{sec:architecture}

\subsection{Project Structure}

The project is organized into a modular directory tree, separating agent definitions, systems (spatial index, glucose field, DNA, effects), treatment logic, parameter management, the model driver, data collection, machine learning, visualization, and tests.

\begin{lstlisting}[language={}, caption={Project directory tree.}, label={lst:project-tree}, basicstyle=\ttfamily\footnotesize, numbers=none, frame=single]
rcc_repast4py/
|-- run.py                        # Entry point
|-- config/
|   `-- default_params.yaml       # Default parameter file
|-- src/
|   |-- agents/
|   |   |-- agent_types.py        # AgentType IntEnum (0-17)
|   |   |-- cell.py               # Base Cell(ABC, Agent) class
|   |   |-- tumor_cell.py         # TumorCell with Warburg glucose
|   |   |-- cd8_cytotoxic_t_cell.py
|   |   |-- cd8_t_cell.py
|   |   |-- cd4_t_cell.py
|   |   |-- cd4_t_cell_h1.py
|   |   |-- cd4_t_cell_h2.py
|   |   |-- regulatory_t_cell.py
|   |   |-- dendritic_cell.py
|   |   |-- plasmacitoid_dc.py
|   |   |-- macrophage_m1.py
|   |   |-- macrophage_m2.py
|   |   |-- natural_killer.py
|   |   |-- mast_cell.py
|   |   |-- neutrophil.py
|   |   |-- adipocyte.py
|   |   |-- blood.py
|   |   |-- sex_hormone.py
|   |   |-- cytokine.py
|   |   `-- t_cell.py             # Shared T-cell base
|   |-- systems/
|   |   |-- grid_utils.py         # 3D neighbor queries
|   |   |-- glucose_field.py      # GlucoseField class
|   |   |-- dna.py                # 16-gene DNA mutation system
|   |   |-- effect.py             # Effect accumulator
|   |   `-- hormonal_effect.py    # Hormonal effect system
|   |-- treatments/
|   |   |-- drug.py               # Drug base class
|   |   |-- ici.py                # ICIDrug
|   |   |-- tki.py                # TKIDrug
|   |   `-- treatment.py          # Treatment container
|   |-- parameters/
|   |   |-- parameters.py         # Base Parameters dataclass
|   |   |-- model_parameters.py   # ModelParameters (4 params)
|   |   |-- patient_parameters.py # PatientParameters (18 params)
|   |   `-- weight_parameters.py  # WeightParameters (92 params)
|   |-- model/
|   |   |-- rcc_model.py          # RCCModel main driver
|   |   |-- cell_adder.py         # Agent initialization
|   |   |-- observer.py           # Kill/event counters
|   |   `-- measures_utils.py     # Volume/concentration helpers
|   |-- data/
|   |   `-- data_collector.py     # Data collection utilities
|   |-- learning/
|   |   |-- train_model.py        # Optuna optimization
|   |   |-- best_parameters.py    # Best parameter storage
|   |   `-- data/
|   |       `-- ARON.csv          # ARON-1 clinical dataset
|   `-- visualization/
|       |-- plot_results.py       # Population dynamics plots
|       `-- plot_glucose.py       # Glucose field visualization
`-- tests/
    |-- test_dna.py
    |-- test_effects.py
    |-- test_glucose_field.py
    `-- test_single_step.py
\end{lstlisting}

\subsection{Plain Python Model Class}

Unlike Mesa, which requires the model class to inherit from \texttt{mesa.Model} and use Mesa's scheduler abstraction, the Repast4Py implementation uses a plain Python class, \texttt{RCCModel}. This class does not inherit from any framework base class; instead, it manually manages:

\begin{itemize}
    \item A \texttt{SharedContext} instance for agent registration with Repast4Py.
    \item A \texttt{SharedGrid} instance for spatial placement.
    \item A custom \texttt{spatial\_index} dictionary for efficient neighbor queries.
    \item A \texttt{GlucoseField} instance for metabolic dynamics.
    \item Agent lifecycle methods: \texttt{add\_agent()}, \texttt{remove\_agent()}, \texttt{move\_agent()}.
    \item A standard library \texttt{random.Random} instance (\texttt{self.rng}) seeded deterministically.
\end{itemize}

There is no framework-imposed scheduler. Instead, Repast4Py's \texttt{schedule.init\_schedule\_runner()} creates an event-driven runner that invokes \texttt{model.step()} at each tick. The \texttt{step()} method itself handles the ordering of sub-steps (effect collection, treatment application, glucose field update, agent stepping, deferred movement, and hormone spawning) explicitly, providing full control over the simulation loop.

Figure~\ref{fig:model_architecture} illustrates the overall architecture of the RCCModel, showing the main components and their data flow relationships.

\begin{figure}[H]
\centering
% Model Architecture — component diagram
\begin{tikzpicture}[
    main/.style={rectangle, rounded corners=4pt, draw, thick, fill=primary!15,
                 font=\small\bfseries, minimum width=2.8cm, minimum height=1.2cm,
                 text centered},
    comp/.style={rectangle, rounded corners=3pt, draw, fill=secondary!12,
                 font=\small, minimum width=2.4cm, minimum height=0.9cm,
                 align=center, inner sep=3pt},
    arrow/.style={->, >=stealth, thick}
]

% Central node
\node[main] (model) at (5,4) {RCCModel};

% Left — Agent Hierarchy
\node[comp] (agents) at (0.5,4) {Agent Hierarchy\\ 18 types};
\draw[arrow] (model) -- (agents) node[midway, above, font=\footnotesize] {manages};

% Top — Glucose Field
\node[comp] (glucose) at (5,7) {Glucose Field\\ 3D diffusion};
\draw[arrow] (model) -- (glucose) node[midway, right, font=\footnotesize] {updates};

% Right — Effect System
\node[comp] (effects) at (9.5,5.5) {Effect System\\ 14 paracrine fields};
\draw[arrow] (model) -- (effects) node[midway, above, font=\footnotesize, sloped] {collects};

% Bottom-right — Treatment
\node[comp] (treat) at (9.5,2.5) {Treatment\\ ICI / TKI};
\draw[arrow] (model) -- (treat) node[midway, below, font=\footnotesize, sloped] {applies};

% Bottom — Parameter Management
\node[comp] (params) at (5,1) {Parameters\\ YAML config};
\draw[arrow] (params) -- (model) node[midway, right, font=\footnotesize] {configures};

% Bottom-left — Data Collection
\node[comp] (data) at (0.5,1.5) {Data Collection\\ CSV logging};
\draw[arrow] (model) -- (data) node[midway, left, font=\footnotesize] {records};

% Cross-component flows (dashed)
\draw[arrow, dashed, neutral] (effects) -- (agents)
    node[midway, above, font=\footnotesize] {modifies};
\draw[arrow, dashed, neutral] (glucose) -| (agents)
    node[pos=0.25, above, font=\footnotesize] {feeds};
\draw[arrow, dashed, neutral] (treat) -- (effects)
    node[midway, right, font=\footnotesize] {alters};

\end{tikzpicture}
\caption{Architecture of the RCCModel system. The main model coordinates interactions between the glucose field (3D diffusion), agent management (cell adder, spatial index, shared grid), data collection (observer), and treatment systems. The step sequence annotation shows the ordered execution of simulation phases.}
\label{fig:model_architecture}
\end{figure}

% ======================================================================
\section{Spatial Index Design}
\label{sec:spatial-index}

\subsection{Problem}

Repast4Py's \texttt{SharedGrid} provides basic spatial placement and movement of agents, but it does \emph{not} expose the \texttt{get\_neighbors()} or \texttt{get\_neighborhood()} methods that Mesa's grids provide. This means there is no built-in mechanism to query which agents are adjacent to a given position---a fundamental operation in the RCC model, where every agent interaction (phagocytosis, antigen presentation, killing, effect propagation) depends on spatial proximity.

\subsection{Solution}

The model maintains a manual spatial index implemented as a \texttt{defaultdict(list)} that maps each \texttt{(x, y, z)} coordinate tuple to the list of agents currently occupying that position:

\begin{lstlisting}[caption={Spatial index declaration in \texttt{RCCModel.\_\_init\_\_()}.}, label={lst:spatial-index}]
from collections import defaultdict

self.spatial_index = defaultdict(list)
\end{lstlisting}

This index is updated on every agent lifecycle operation:
\begin{itemize}
    \item \texttt{add\_agent(agent)}: appends the agent to \texttt{spatial\_index[agent.pos]}.
    \item \texttt{remove\_agent(agent)}: removes the agent from \texttt{spatial\_index[agent.pos]}.
    \item \texttt{move\_agent(agent, new\_pos)}: removes from the old position list and appends to the new one.
\end{itemize}

\subsection{Grid Utilities Module}

The \texttt{src/systems/grid\_utils.py} module provides five functions that operate on the spatial index:

\begin{enumerate}
    \item \texttt{get\_neighborhood\_3d(width, height, depth, pos, radius, moore)}: Returns all coordinate tuples in the Moore (Chebyshev) or Von Neumann (Manhattan) neighborhood of a given position, excluding the center. For Moore neighborhoods at radius~1 in 3D, this yields up to 26 neighbors ($3^3 - 1$).

    \item \texttt{get\_agents\_at\_3d(spatial\_index, pos)}: Returns the list of agents at a specific coordinate.

    \item \texttt{get\_neighbors\_3d(spatial\_index, grid\_dims, pos, radius, moore)}: Returns all agents in the neighborhood of a position by iterating over the neighborhood coordinates and collecting agents from the spatial index.

    \item \texttt{is\_cell\_empty\_3d(spatial\_index, pos)}: Returns \texttt{True} if no agents occupy the given position.

    \item \texttt{get\_empty\_neighbors\_3d(spatial\_index, grid\_dims, pos, radius, moore)}: Returns the list of neighboring positions that have no agents, used for cell duplication and random movement.
\end{enumerate}

\begin{lstlisting}[caption={Moore neighborhood computation in 3D.}, label={lst:neighborhood-3d}]
def get_neighborhood_3d(width, height, depth, pos, radius=1, moore=True):
    x, y, z = pos
    neighborhood = []
    for dx in range(-radius, radius + 1):
        for dy in range(-radius, radius + 1):
            for dz in range(-radius, radius + 1):
                if dx == 0 and dy == 0 and dz == 0:
                    continue
                if not moore and abs(dx) + abs(dy) + abs(dz) > radius:
                    continue
                nx, ny, nz = x + dx, y + dy, z + dz
                if 0 <= nx < width and 0 <= ny < height and 0 <= nz < depth:
                    neighborhood.append((nx, ny, nz))
    return neighborhood
\end{lstlisting}

The boundary check (\texttt{0 <= nx < width}, etc.) ensures that the neighborhood is clipped at grid edges, consistent with the sticky border semantics of the \texttt{SharedGrid}.

% ======================================================================
\section{Agent Serialization}
\label{sec:agent-serialization}

\subsection{Agent Identity in Repast4Py}

Repast4Py requires every agent to be constructed with three arguments: an integer \texttt{id}, an integer \texttt{type}, and an integer \texttt{rank} (the MPI rank that owns the agent). The framework stores these as a \texttt{uid} tuple \texttt{(id, type, rank)} on every agent. This design supports distributed serialization, as agents can be uniquely identified and routed across MPI ranks.

\subsection{Agent Type Registry}

The model defines an \texttt{AgentType} \texttt{IntEnum} that maps all 18 agent types to integer IDs in the range $[0, 17]$:

\begin{lstlisting}[caption={Agent type registry (\texttt{src/agents/agent\_types.py}).}, label={lst:agent-types}]
from enum import IntEnum

class AgentType(IntEnum):
    TUMOR_CELL = 0
    CD8_CYTOTOXIC_T_CELL = 1
    CD8_NAIVE_T_CELL = 2
    CD4_NAIVE_T_CELL = 3
    CD4_HELPER1_T_CELL = 4
    CD4_HELPER2_T_CELL = 5
    REGULATORY_T_CELL = 6
    DENDRITIC_CELL = 7
    PLASMACITOID_DC = 8
    MACROPHAGE_M1 = 9
    MACROPHAGE_M2 = 10
    NATURAL_KILLER = 11
    MAST_CELL = 12
    NEUTROPHIL = 13
    ADIPOCYTE = 14
    BLOOD = 15
    SEX_HORMONE = 16
    CYTOKINE = 17
\end{lstlisting}

\subsection{The \texttt{save()} Method}

Each agent implements a \texttt{save()} method that returns a tuple of \texttt{(uid, state\_tuple)} for distributed serialization. The base implementation in \texttt{Cell} is:

\begin{lstlisting}[caption={Agent serialization via \texttt{save()}.}, label={lst:save-method}]
def save(self):
    return (self.uid, (self.pos, self._alive))
\end{lstlisting}

Subclasses extend the state tuple to include type-specific fields (e.g., the \texttt{TumorCell} includes its \texttt{DNA} state, blood access flag, and experienced effects).

\subsection{Type Identification}

Because Repast4Py agents carry their type as the second element of the \texttt{uid} tuple, type identification uses integer comparison rather than Python's \texttt{isinstance()} check:

\begin{lstlisting}[caption={Type identification pattern.}, label={lst:type-id}]
# Instead of isinstance(agent, TumorCell):
if agent.uid[1] == AgentType.TUMOR_CELL:
    # handle tumor cell
\end{lstlisting}

This approach is both faster (integer comparison vs.\ method resolution order traversal) and compatible with ghost agents that may not carry the full Python class hierarchy across MPI ranks.

% ======================================================================
\section{Mesa to Repast4Py Porting}
\label{sec:porting}

The original RCC model was implemented in Mesa~\cite{mesa, masl_paper_a}. Porting it to Repast4Py required systematic translation of Mesa idioms to their Repast4Py equivalents. Table~\ref{tab:porting-mapping} summarizes the key mappings.

\begin{table}[H]
\centering
\caption{Key API mapping from Mesa to Repast4Py.}
\label{tab:porting-mapping}
\begin{tabular}{@{}p{6cm}p{7.5cm}@{}}
\toprule
\textbf{Mesa} & \textbf{Repast4Py} \\
\midrule
\texttt{mesa.Agent(model)} & \texttt{repast4py.core.Agent(id, type, rank)} \\
\texttt{self.model.random} & \texttt{model.rng} (stdlib \texttt{random.Random}) \\
\texttt{self.model.grid.get\_neighbors()} & \texttt{grid\_utils.get\_neighbors\_3d(spatial\_index, ...)} \\
\texttt{isinstance(agent, TumorCell)} & \texttt{agent.uid[1] == AgentType.TUMOR\_CELL} \\
\texttt{self.model.schedule.agents} & \texttt{model.\_all\_agents} \\
\texttt{model.agents\_by\_type[Type]} & \texttt{model.get\_agents\_by\_type\_id(type\_id)} \\
\texttt{TumorCell.blood\_source} (class var) & \texttt{model.tumor\_blood\_sources} \\
\texttt{Cell.search\_dimension} (class var) & \texttt{model.cell\_search\_dimension} \\
\bottomrule
\end{tabular}
\end{table}

\subsection{Random Number Generation}

Mesa provides a model-level \texttt{self.random} attribute backed by a NumPy random generator. In Repast4Py, no such attribute is provided, so the model creates a standard library \texttt{random.Random} instance seeded with \texttt{random\_seed + rank}:

\begin{lstlisting}[caption={RNG initialization in RCCModel.}]
import random
self.rng = random.Random(self.model_params.random_seed + self.rank)
\end{lstlisting}

All agent behaviors reference \texttt{self.model.rng} for stochastic decisions, ensuring reproducibility at a given seed.

\subsection{Class-Level State to Model-Level State}

In the Mesa version, certain state was stored as class variables (e.g., \texttt{TumorCell.blood\_source} tracked all tumor cells with blood vessel access). In the Repast4Py version, this state is moved to the model instance to avoid issues with class-level state in a distributed setting:

\begin{lstlisting}[caption={Class-level state moved to model.}]
# In RCCModel.__init__():
self.tumor_blood_sources = []
self.cell_search_dimension = 5
\end{lstlisting}

\subsection{Deferred Movement}

The deferred movement pattern from the Mesa version is preserved. During their \texttt{step()} method, agents set a \texttt{\_desired\_pos} attribute rather than moving immediately. After all agents have been stepped, the model calls \texttt{move\_all\_agents()}, which resolves all pending movements in a single pass:

\begin{lstlisting}[caption={Deferred movement resolution.}, label={lst:deferred-move}]
def move_all_agents(self):
    for agent in list(self._all_agents):
        if (getattr(agent, '_alive', True)
                and isinstance(agent, Cell)
                and hasattr(agent, '_desired_pos')
                and agent._desired_pos is not None
                and agent._desired_pos != agent.pos):
            self.move_agent(agent, agent._desired_pos)
\end{lstlisting}

This approach prevents order-of-execution artifacts where early-stepping agents could influence the spatial configuration seen by later-stepping agents within the same tick.

% ======================================================================
\section{Glucose Field Implementation}
\label{sec:glucose-field-impl}

The glucose concentration field is implemented as a standalone class, \texttt{GlucoseField}, in \texttt{src/systems/glucose\_field.py}. It uses pure NumPy arrays and SciPy convolution rather than Repast4Py's built-in value layers, for two reasons: (i) full control over the diffusion kernel and boundary conditions, and (ii) performance, since NumPy operations are vectorized at the C level.

\subsection{Data Structure}

The field is stored as a NumPy 3D array of shape \texttt{(width, height, depth)} with \texttt{float64} precision:

\begin{lstlisting}[caption={Glucose field initialization.}]
self.field = np.full((width, height, depth), initial_concentration,
                     dtype=np.float64)
\end{lstlisting}

\subsection{Diffusion Kernel}

Diffusion is computed via the 3D discrete Laplacian, applied using \texttt{scipy.ndimage.convolve}. The kernel is a $3 \times 3 \times 3$ array with:
\begin{itemize}
    \item Center value: $-6$
    \item Six face-adjacent neighbors (along $\pm x$, $\pm y$, $\pm z$): $+1$ each
    \item All other entries: $0$ (edge and corner neighbors are not included in the 6-connected Laplacian)
\end{itemize}

\begin{lstlisting}[caption={Laplacian kernel construction.}]
self._kernel = np.zeros((3, 3, 3), dtype=np.float64)
self._kernel[1, 1, 0] = 1.0   # z-1
self._kernel[1, 1, 2] = 1.0   # z+1
self._kernel[1, 0, 1] = 1.0   # y-1
self._kernel[1, 2, 1] = 1.0   # y+1
self._kernel[0, 1, 1] = 1.0   # x-1
self._kernel[2, 1, 1] = 1.0   # x+1
self._kernel[1, 1, 1] = -6.0  # center
\end{lstlisting}

\subsection{Boundary Conditions}

Zero-flux (Neumann) boundary conditions are enforced by passing \texttt{mode='nearest'} to the convolution function. This causes boundary voxels to replicate their own values when the kernel extends beyond the grid, effectively setting the normal derivative of the concentration to zero at all boundaries. The gradient computation uses clamped central finite differences with the same Neumann semantics.

\subsection{Separation from Repast4Py Value Layers}

Repast4Py provides its own \texttt{SharedValueLayer} for continuous fields. However, the glucose field implementation deliberately avoids this abstraction for several reasons:
\begin{enumerate}
    \item The \texttt{SharedValueLayer} is designed for distributed access patterns that introduce synchronization overhead unnecessary for single-rank execution.
    \item Direct NumPy array access allows use of SciPy's optimized convolution routines without intermediate copying.
    \item The pure-NumPy approach simplifies debugging and testing, as the field can be inspected and manipulated with standard array operations.
\end{enumerate}

% ======================================================================
\section{Entry Point and Execution}
\label{sec:entry-point}

The simulation is launched from \texttt{run.py}, which follows the standard Repast4Py execution pattern:

\begin{lstlisting}[caption={Simulation entry point (\texttt{run.py}).}, label={lst:run-py}]
from mpi4py import MPI
from repast4py import schedule
from src.model.rcc_model import RCCModel

def main():
    comm = MPI.COMM_WORLD
    rank = comm.Get_rank()

    model = RCCModel(comm=comm)

    runner = schedule.init_schedule_runner(comm)
    runner.schedule_repeating_event(1, 1, model.step)
    runner.schedule_stop(model.max_steps)
    runner.schedule_end_event(model.log_data)

    runner.execute()

if __name__ == '__main__':
    main()
\end{lstlisting}

\subsection{Execution Command}

The simulation is invoked via MPI:

\begin{lstlisting}[language=bash, numbers=none]
mpirun -n 1 python run.py
\end{lstlisting}

The \texttt{-n 1} flag specifies a single MPI rank, which is the current target configuration. The architecture supports multi-rank execution by increasing the rank count; in that scenario, the \texttt{SharedContext} would distribute agents across ranks and the \texttt{SharedGrid} would handle ghost-agent synchronization via buffer regions.

\subsection{Simulation Loop}

The Repast4Py schedule runner calls \texttt{model.step()} at every tick (starting at tick 1, repeating every 1 tick). Within \texttt{step()}, the following sub-steps are executed in order:

\begin{enumerate}
    \item \textbf{Data collection}: Record current population counts and glucose field statistics.
    \item \textbf{Termination check}: Evaluate whether the simulation should stop (tumor eliminated, maximum tumor count reached, growth threshold exceeded, or maximum steps reached).
    \item \textbf{Effect propagation}: Cells collect and apply effects from their neighbors.
    \item \textbf{Treatment}: If the current step exceeds the treatment start time, apply ICI and/or TKI drug effects.
    \item \textbf{Sex drift}: Apply sex-specific CD8+ T cell population adjustment.
    \item \textbf{Glucose field step}: Inject glucose from blood vessels, diffuse, and decay.
    \item \textbf{Blood management}: Angiogenesis and blood access refresh for tumor cells.
    \item \textbf{Search dimension update}: Adjust the agent search radius based on tumor population size.
    \item \textbf{Agent stepping}: Shuffle all agents and execute their individual \texttt{step()} methods.
    \item \textbf{Deferred movement}: Resolve all pending agent movements.
    \item \textbf{Hormone spawning}: Spawn sex hormones from blood vessels.
    \item \textbf{Step counter increment}.
\end{enumerate}

When the simulation terminates (either by the schedule runner reaching \texttt{max\_steps} or by the model setting \texttt{self.running = False}), the end event triggers \texttt{model.log\_data()}, which writes the collected time-series data to \texttt{logs/simulation\_log.csv}.

Figure~\ref{fig:simulation_pipeline} provides a detailed view of the simulation pipeline, showing the complete sequence of operations executed in each time step and the decision points that control simulation flow.

\begin{figure}[H]
\centering
\begin{tikzpicture}[
  box/.style={rectangle, rounded corners, minimum width=2.8cm, minimum height=1cm, text centered, font=\footnotesize, draw=black, thick},
  databox/.style={box, fill=blue!15},
  processbox/.style={box, fill=green!15},
  agentbox/.style={box, fill=orange!15},
  controlbox/.style={box, fill=red!15},
  arrow/.style={->, >=stealth, thick, shorten >=2pt, shorten <=2pt},
  decision/.style={diamond, aspect=2, minimum width=2cm, minimum height=0.8cm, text centered, font=\tiny, draw=black, thick, fill=yellow!15}
]

% Main pipeline flow (vertical)
\node[databox] (collect) at (0, 8) {\textbf{1. Collect Data}\\Population counts\\Glucose statistics};

\node[decision] (terminate) at (0, 6.5) {\textbf{2. Check}\\Termination\\Conditions};

\node[processbox] (effects) at (0, 5) {\textbf{3. Apply Effects}\\Collect \& sum\\neighbor effects};

\node[controlbox] (treatment) at (0, 3.5) {\textbf{4. Treatment}\\ICI \& TKI\\drug effects};

\node[controlbox] (sexdrift) at (0, 2) {\textbf{5. Sex Drift}\\CD8+ modulation\\by sex};

\node[processbox] (glucose) at (0, 0.5) {\textbf{6. Glucose Step}\\Diffusion \& decay\\Blood injection};

\node[processbox] (blood) at (0, -1) {\textbf{7. Manage Blood}\\Angiogenesis\\BFS refresh};

\node[processbox] (search) at (0, -2.5) {\textbf{8. Search Dimension}\\Update cell\\search radius};

\node[agentbox] (agents) at (0, -4) {\textbf{9. Agent Steps}\\Individual agent\\behaviors (shuffled)};

\node[agentbox] (move) at (0, -5.5) {\textbf{10. Deferred Move}\\Batch movement\\Update positions};

\node[controlbox] (hormones) at (0, -7) {\textbf{11. Spawn Hormones}\\Sex hormone\\generation};

% Terminal paths
\node[controlbox] (log) at (4, 6.5) {\textbf{Write Log}\\\textbf{Terminate}};
\node[controlbox] (continue) at (-4, 5) {\textbf{Next Step}\\$t \leftarrow t+1$};

% Main flow arrows
\draw[arrow] (collect) -- (terminate);
\draw[arrow] (terminate) -- (effects);
\draw[arrow] (effects) -- (treatment);
\draw[arrow] (treatment) -- (sexdrift);
\draw[arrow] (sexdrift) -- (glucose);
\draw[arrow] (glucose) -- (blood);
\draw[arrow] (blood) -- (search);
\draw[arrow] (search) -- (agents);
\draw[arrow] (agents) -- (move);
\draw[arrow] (move) -- (hormones);

% Termination decision
\draw[arrow] (terminate) -- (log) node[midway, above] {\scriptsize yes};
\draw[arrow] (terminate) -- (effects) node[midway, right] {\scriptsize no};

% Loop back
\draw[arrow] (hormones) to [out=180, in=270] (continue);
\draw[arrow] (continue) to [out=90, in=180] (collect);

% Side annotations
\begin{scope}[xshift=6cm]
% Data collection details
\node[text width=3cm, font=\tiny, fill=gray!5, rounded corners, draw] at (0, 8) {
\textbf{Data collected:}\\
• Agent counts by type\\
• Glucose: mean, min, max\\
• Kill statistics\\
• Growth rates
};

% Termination conditions
\node[text width=3cm, font=\tiny, fill=gray!5, rounded corners, draw] at (0, 6) {
\textbf{Terminal conditions:}\\
• Max steps reached\\
• Tumor eliminated\\
• Tumor $>$ threshold\\
• Growth rate $>$ limit
};

% Agent behaviors
\node[text width=3cm, font=\tiny, fill=gray!5, rounded corners, draw] at (0, -4) {
\textbf{Agent behaviors:}\\
• Move towards targets\\
• Kill/proliferate\\
• Differentiate\\
• Apply hormones\\
• Set desired position
};
\end{scope}

% Parallel processes annotation
\begin{scope}[xshift=-6cm]
\node[text width=3cm, font=\tiny, fill=gray!5, rounded corners, draw] at (0, 0) {
\textbf{Glucose field:}\\
• 3D diffusion (PDE)\\
• Blood injection\\
• Cell consumption\\
• Natural decay\\
• Gradient computation
};

\node[text width=3cm, font=\tiny, fill=gray!5, rounded corners, draw] at (0, -2.5) {
\textbf{Spatial management:}\\
• Moore neighborhoods\\
• 3D distance queries\\
• BFS from blood vessels\\
• Deferred movement\\
• Collision handling
};
\end{scope}

% Time annotation
\node[text width=8cm, font=\footnotesize, text centered] at (0, -8.5) {
\textbf{One Simulation Step} --- Process repeats until termination condition met
};

\end{tikzpicture}
\caption{Simulation pipeline for the RCC model. Each time step follows a fixed sequence of 11 operations, from data collection through hormone spawning. The termination condition check determines whether to continue the simulation or write final results. Side annotations detail the specific operations performed in each phase.}
\label{fig:simulation_pipeline}
\end{figure}
